\section{Evaluation metrics}

\subsection{PSDS1}
PSDS1 is threshold-independent. For class $c$ and an operating point,
\begin{equation}
r_c=\frac{TP_c}{TP_c+FN_c},
\qquad
f_c=\frac{FP_c\cdot 3600}{T_{\mathrm{ds}}},
\end{equation}
where $T_{\mathrm{ds}}$ is the evaluation duration in seconds and $f_c$ is false
positives per hour. A prediction matches a ground-truth event of the same class
when their temporal IoU is at least $0.5$. Published DCASE PSDS1 uses the
DTC/GTC criterion (typically $0.7$); this report follows the paper's IoU-0.5
wording and labels the resulting score PSDS1@IoU0.5 where the distinction
matters. Per-class curves are averaged:
\begin{equation}
\bar r(f)=\frac{1}{C}\sum_{c=1}^{C}r_c(f),
\qquad
\mathrm{PSDS1}
=
\frac{1}{e_{\max}}
\int_0^{e_{\max}}\bar r(f)\,df,
\end{equation}
with $e_{\max}=100$ false positives per hour. Following the paper, PSDS1 sets
the inter-class variance penalty to zero ($\alpha_{\mathrm{ST}}=0$); PSDS2 uses
$\alpha_{\mathrm{ST}}=1$.

\subsection{PSDS2}
PSDS2 is the DCASE scenario-2 metric as implemented here: DTC $=0.1$,
GTC $=0.1$, CTTC $=0.3$ (this is our configuration, not a metric reported by the
source paper),
cross-trigger weight $\alpha_{\mathrm{CT}}=0.5$, and variance penalty
$\alpha_{\mathrm{ST}}=1$. The effective false-positive rate adds cross-triggers:
\begin{equation}
e_c
=
f_c
+
\alpha_{\mathrm{CT}}
\frac{1}{C-1}
\sum_{k\ne c}
\frac{CT_{c,k}\cdot 3600}{T_k},
\end{equation}
where $CT_{c,k}$ is the number of class-$c$ detections cross-triggering class
$k$, and $T_k$ is the total activity duration of class $k$. The effective TPR
penalises class instability:
\begin{equation}
\tilde r(f)=\bar r(f)-\alpha_{\mathrm{ST}}\,\sigma_r(f),
\end{equation}
where $\sigma_r(f)$ is the standard deviation of the per-class TPRs at
effective FPR $f$.

\subsection{Collar-F1}
Collar-F1 matches an event when
\begin{equation}
|\hat t_{\on}-t_{\on}|\le 0.2,
\qquad
|\hat t_{\off}-t_{\off}|\le \max(0.2,\;0.2\,d),
\end{equation}
where $d$ is the ground-truth duration. With matched counts,
\begin{equation}
P=\frac{TP}{TP+FP},
\qquad
R=\frac{TP}{TP+FN},
\qquad
F_1=\frac{2PR}{P+R}.
\end{equation}
Collar-F1 is threshold-dependent; PSDS1 and PSDS2 are not.
